\documentclass[10pt,a4paper]{article}
\usepackage[utf8]{inputenc}
\usepackage[vietnamese]{babel}
\usepackage{amsmath,amssymb,amsfonts}
\usepackage{geometry}
\geometry{a4paper, left=10mm, right=10mm, top=10mm, bottom=10mm}
\usepackage{tikz}
\usetikzlibrary{shapes,arrows.meta,positioning,calc}
\usepackage[most]{tcolorbox}
\usepackage{array,booktabs,multicol}
\usepackage{fontawesome}
\usepackage{helvet}
\renewcommand{\familydefault}{\sfdefault}
\usepackage{parskip}

\definecolor{primary}{RGB}{112, 48, 160}
\definecolor{teal}{RGB}{0, 128, 128}
\definecolor{orange}{RGB}{230, 108, 44}
\definecolor{green}{RGB}{46, 139, 87}
\definecolor{rose}{RGB}{192, 0, 0}

\pagestyle{empty}

\begin{document}

% --- HEADER BANNER ---
\begin{tcolorbox}[
  enhanced,
  colback=violet!8!white,
  colframe=violet!80!black,
  arc=2.5mm, boxrule=1pt,
  left=3mm, right=3mm, top=2mm, bottom=2mm
]
  \begin{tabular*}{\textwidth}{@{\extracolsep{\fill}} lr}
    {\small\bfseries\color{violet!85!black}\faBook\ SỔ TAY KIẾN THỨC TOÁN 2} & 
    {\small\bfseries\color{teal!85!black}\faGraduationCap\ CHƯƠNG TRÌNH GDPT 2018}
  \end{tabular*}
  \vspace{-1mm}
  \begin{center}
    \rule{0.98\textwidth}{0.6pt}
  \end{center}
  \vspace{-2mm}
  \centering
  {\LARGE\bfseries\color{violet!90!black}\faBookmark\ BÀI 50: SO SÁNH CÁC SỐ TRÒN TRĂM, TRÒN CHỤC}\\[1mm]
  {\itshape\footnotesize\color{black!75} Quy tắc so sánh số tròn trăm, số tròn chục và thứ tự các số trên 1 trang A4}
\end{tcolorbox}

\vspace{-2mm}

\noindent
\begin{minipage}[t]{0.49\textwidth}\vspace{0pt}
\begin{tcolorbox}[
  enhanced,
  colback=violet!4!white,
  colframe=violet!75!black,
  arc=2mm, boxrule=0.8pt,
  left=2mm, right=2mm, top=1.8mm, bottom=1.8mm,
  title={\footnotesize\bfseries\color{white}\faKey\ 1. So sánh các số tròn trăm}
]
{\footnotesize
\begin{itemize}\setlength{\itemsep}{1pt}\setlength{\leftskip}{-2.5mm}
  \item Để so sánh hai số tròn trăm, ta \textbf{so sánh chữ số hàng trăm}:
  \begin{itemize}\setlength{\itemsep}{0.5pt}\setlength{\leftskip}{-2mm}
    \item Số nào có chữ số hàng trăm lớn hơn thì \textbf{lớn hơn}.
    \item Số nào có chữ số hàng trăm bé hơn thì \textbf{bé hơn}.
  \end{itemize}
  \item \textbf{Ví dụ:} So sánh $300$ và $500$:
  \begin{itemize}\setlength{\itemsep}{0.5pt}\setlength{\leftskip}{-2mm}
    \item Chữ số hàng trăm là $3$ và $5$.
    \item Vì $3 < 5$ nên $\mathbf{300 < 500}$ (hay $500 > 300$).
  \end{itemize}
  \item Đặc biệt: $1\,000$ có $4$ chữ số nên lớn hơn mọi số tròn trăm có $3$ chữ số ($1\,000 > 900$).
\end{itemize}
}
\end{tcolorbox}
\end{minipage}
\hfill
\begin{minipage}[t]{0.49\textwidth}\vspace{0pt}
\begin{tcolorbox}[
  enhanced,
  colback=teal!4!white,
  colframe=teal!75!black,
  arc=2mm, boxrule=0.8pt,
  left=2mm, right=2mm, top=1.8mm, bottom=1.8mm,
  title={\footnotesize\bfseries\color{white}\faSliders\ 2. So sánh các số tròn chục}
]
{\footnotesize
\begin{itemize}\setlength{\itemsep}{1pt}\setlength{\leftskip}{-2.5mm}
  \item Khi so sánh hai số tròn chục có 3 chữ số:
  \begin{itemize}\setlength{\itemsep}{0.5pt}\setlength{\leftskip}{-2mm}
    \item \textbf{Bước 1:} So sánh chữ số \textbf{hàng trăm}.
    \item \textbf{Bước 2:} Nếu hàng trăm bằng nhau, so sánh chữ số \textbf{hàng chục}.
  \end{itemize}
  \item \textbf{Ví dụ 1:} So sánh $240$ và $180$:
  \begin{itemize}\setlength{\itemsep}{0.5pt}\setlength{\leftskip}{-2mm}
    \item Hàng trăm: $2 > 1 \implies \mathbf{240 > 180}$.
  \end{itemize}
  \item \textbf{Ví dụ 2:} So sánh $430$ và $470$:
  \begin{itemize}\setlength{\itemsep}{0.5pt}\setlength{\leftskip}{-2mm}
    \item Hàng trăm: $4 = 4$. Hàng chục: $3 < 7$.
    \item Do đó: $\mathbf{430 < 470}$.
  \end{itemize}
\end{itemize}
}
\end{tcolorbox}
\end{minipage}

\vspace{-1.5mm}
\begin{tcolorbox}[
  enhanced,
  colback=orange!4!white,
  colframe=orange!80!black,
  arc=2mm, boxrule=0.8pt,
  left=2mm, right=2mm, top=1.5mm, bottom=1.5mm,
  title={\footnotesize\bfseries\color{white}\faSitemap\ 3. Sơ đồ tư duy: Các bước so sánh số tròn trăm, tròn chục}
]
\centering
\begin{tikzpicture}[x=1.35cm, y=0.8cm]
  \tikzset{arr/.style={-{Stealth[scale=0.8]}}, every node/.append style={align=center}}
  \node[rectangle, rounded corners=2pt, draw=violet!80!black, fill=violet!15, font=\scriptsize\bfseries] (start) at (0,0) {SO SÁNH\\HAI SỐ};
  \node[rectangle, rounded corners=2pt, draw=teal!80!black, fill=teal!15, font=\scriptsize\bfseries] (step1) at (2.4,0) {So sánh\\HÀNG TRĂM};
  \node[rectangle, rounded corners=2pt, draw=orange!80!black, fill=orange!15, font=\scriptsize\bfseries] (diff) at (5.2,0.8) {Khác nhau\\$\implies$ Kết luận luôn};
  \node[rectangle, rounded corners=2pt, draw=green!70!black, fill=green!15, font=\scriptsize\bfseries] (same) at (5.2,-0.8) {Bằng nhau\\$\implies$ So tiếp HÀNG CHỤC};

  \draw[arr, thick] (start) -- (step1);
  \draw[arr, thick] (step1) |- (diff);
  \draw[arr, thick] (step1) |- (same);
\end{tikzpicture}
\end{tcolorbox}

\vspace{-1.5mm}
\noindent
\begin{minipage}[t]{0.49\textwidth}\vspace{0pt}
\begin{tcolorbox}[
  enhanced,
  colback=green!4!white,
  colframe=green!60!black,
  arc=2mm, boxrule=0.8pt,
  left=2mm, right=2mm, top=1.8mm, bottom=1.8mm,
  title={\footnotesize\bfseries\color{white}\faCheckCircle\ 4. Ví dụ mẫu điển hình}
]
{\footnotesize
\textbf{1. Điền dấu ($>, <, =$):}\\
- $600\ \dots\ 400 \implies 6 > 4 \implies 600 > 400$.\\
- $520\ \dots\ 580 \implies 5 = 5,\, 2 < 8 \implies 520 < 580$.\\
- $700\ \dots\ 700 \implies 700 = 700$.\\
\textbf{2. Sắp xếp dãy số:} $300, 800, 100, 500$.\\
- Từ bé đến lớn: $\mathbf{100, 300, 500, 800}$.\\
- Từ lớn đến bé: $\mathbf{800, 500, 300, 100}$.
}
\end{tcolorbox}
\end{minipage}
\hfill
\begin{minipage}[t]{0.49\textwidth}\vspace{0pt}
\begin{tcolorbox}[
  enhanced,
  colback=rose!4!white,
  colframe=rose!70!black,
  arc=2mm, boxrule=0.8pt,
  left=2mm, right=2mm, top=1.8mm, bottom=1.8mm,
  title={\footnotesize\bfseries\color{white}\faExclamationTriangle\ 5. Cạm bẫy \& Mẹo ghi nhớ}
]
{\footnotesize
\begin{itemize}\setlength{\itemsep}{1.5pt}\setlength{\leftskip}{-2.5mm}
  \item \textbf{Bẫy số chữ số:} Số có 4 chữ số luôn lớn hơn số có 3 chữ số: $1\,000 > 990$.
  \item \textbf{Mẹo nhìn từ trái sang phải:} Luôn so sánh chữ số ở hàng cao nhất trước (trăm $\to$ chục $\to$ đơn vị).
  \item \textbf{Dấu mở về phía số lớn:} Mũi nhọn luôn quay về số bé hơn, miệng loe quay về số lớn hơn.
\end{itemize}
}
\end{tcolorbox}
\end{minipage}


\end{document}
